\documentclass[10pt,a4paper]{amsart}
\usepackage[margin=19mm]{geometry}
\usepackage{amsmath,amssymb,mathtools}
\usepackage[scr=rsfs,cal=euler]{mathalfa}
\usepackage{xfrac}
\usepackage[unicode,hidelinks,bookmarks=false]{hyperref}
\newcommand{\ie}{i.e.\ }
\newcommand{\cf}{cf.\ }
\newcommand{\resp}{respectively\ }
\newcommand{\Subsection}{\subsection}
\providecommand{\nobreakdash}{\nobreak\hbox{-}}
\setlength{\emergencystretch}{2em}
\multlinegap0pt
\usepackage{bbm}
\usepackage{enumitem}
\theoremstyle{plain}
\newtheorem{theorem}{Theorem}
\newtheorem{lemma}[theorem]{Lemma}
\newtheorem{proposition}[theorem]{Proposition}
\newtheorem{definition}[theorem]{Definition}
\newtheorem{assumption}{Assumption}
\theoremstyle{remark}
\newtheorem{remark}{Remark}
\renewcommand{\hat}{\widehat}
\renewcommand{\tilde}{\widetilde}
\newcommand{\norm}[1]{\left|\left| #1 \right|\right|}
\DeclareMathOperator*{\argmin}{arg\,min}
\title[Consistency of the dual solution]{Efficient Inference on High-Dimensional Linear Models With Missing Outcomes: consistency of the dual solution to the sparse approximation with an estimated propensity score}
\author{Yikun Zhang, Alexander Giessing, Yen-Chi Chen}
\date{Source: arXiv:2309.06429v4 (CC BY 4.0)}
\begin{document}
\maketitle
\noindent\emph{Source and licence.} Efficient Inference on High-Dimensional Linear Models With Missing Outcomes. Version arXiv:2309.06429v4 (CC BY 4.0), distributed under the Creative Commons Attribution 4.0 International licence, \url{http://creativecommons.org/licenses/by/4.0/}, as recorded in the arXiv licence metadata for this exact version and shown on the abstract page \url{https://arxiv.org/abs/2309.06429v4}. Full licence notice: licenses/arxiv-2309.06429-CC-BY-4.0.txt.\par\smallskip

\noindent\textit{Editorial scope.} Counted unit: the article's Theorem with source label \texttt{thm:dual\_consist2} (source0.tex lines 562--574), the consistency of the dual solution of the debiasing program when the propensity score must itself be estimated, delivered together with the article's complete original proof of it (source0.tex lines 2620--2713, one \texttt{proof} environment of 94 lines whose heading names the counted theorem; the article places this proof in its appendix). No mathematics is written, repaired or re-derived: the counted statement and its proof are the article's own text line for line, and no retained line is substituted, re-flowed or omitted. Required context retained uncounted, in source order: the five assumptions the counted statement invokes, namely \texttt{assump:basic} (lines 188--201), \texttt{assump:sub\_gaussian} (lines 438--441), \texttt{assump:gram\_lower\_bnd} (lines 459--463), \texttt{assump:prop\_consistent\_rate} (lines 477--480) and \texttt{assump:sparse\_dual} (lines 522--526); the dual form of the debiasing program (lines 317--320); the definition of the population dual solution and of its sparse approximation (lines 505--513); the statement of the article's known-propensity-score lemma \texttt{lem:dual\_consist} (lines 2513--2521), to which the counted proof reduces its own inequality; the single labelled display \texttt{dual\_taylor} of that lemma's proof (lines 2527--2534), which the counted proof cites by equation number; and the statements (without their proofs) of the five further results the counted proof invokes, namely \texttt{lem:gram\_mat\_conc} (lines 3581--3602), \texttt{lem:dual\_approx\_cone} (lines 3664--3670), \texttt{lem:res\_cross\_poly} (lines 3743--3754), \texttt{lem:cone\_ball\_bnd} (lines 3796--3814) and \texttt{lem:rel\_poly\_cost} (lines 3949--3961). Nothing else of the article is retained: its other results and proofs, its simulation study, its data applications and its remaining appendices are omitted. No citation command occurs in any retained line, so the article's bibliography is not delivered and the wrapper carries no bibliography entry. Every cross-reference inside the retained spans resolves inside the delivered document. Editorial formatting only: an AMS \texttt{amsart} preamble that restates the macros and environments the retained lines use, namely the article's declarations of Theorem, Lemma, Proposition and Definition sharing one counter and of Assumption and Remark with their own counters, the macros \texttt{hat} and \texttt{tilde} renewed as wide accents as in the article, the article's own \texttt{norm} macro and its \texttt{argmin} operator, the \texttt{bbm} package for the indicator symbol and the \texttt{enumitem} package for the labelled enumerations. The article is typeset in its own \texttt{imsart} class with its own style files and fonts; no class or style file is shipped with this package, and the wrapper is an amsart document. Numbering differs from the article, because only a subset of environments is retained: the retained macros are applied to the retained environments in their source order, so the counted theorem prints with the wrapper's own number rather than the article's number, and the retained environments are numbered consecutively by the wrapper. The article's cross-references are not rewritten: the counted proof names the counted theorem and cites the retained displays and the retained statements by their wrapper numbers. No picture environment and no graphical inclusion occurs in any retained line, and the package delivers no image asset. One bundled source file, \texttt{sources/146927/source0.tex}, accompanies the delivered item as the exact source of every retained span. Every retained span is recorded in \texttt{delivery-evidence.json} with method \texttt{exact\_tex}.
\begin{assumption}[High-dimensional sparse linear regression model with MAR outcomes] 
	\leavevmode %Let $(Y, R, X) \in \mathbb{R} \times \{0,1\} \times \mathbb{R}^d$.
	\begin{enumerate}[label=(\alph*)]
		\label{assump:basic}
		\item The data sample $\{(Y_i,R_i,X_i)\}_{i=1}^n$ consists of independent and identically distributed (i.i.d.) observations drawn from $(Y, R, X) \in \mathbb{R} \times \{0,1\} \times \mathbb{R}^d$ generated by the linear model
		\begin{equation}
			\label{linear_reg}
			Y = X^T\beta_0 + \epsilon, \quad \quad \mathrm{E}\left(\epsilon \big| X\right) = 0, \quad\quad \mathrm{E}\left(\epsilon^2 \big| X\right) = \sigma_{\epsilon}^2, 
		\end{equation}
		where $\norm{\beta_0}_0 = \sum_{k=1}^d \mathbbm{1}_{\{\beta_{0k} \neq 0\}} = s_{\beta} \ll d$. 
		
		\item The missingness indicator $R$ is conditionally independent of $Y$ given $X$.
	\end{enumerate}
\end{assumption}

	\begin{equation}
		\label{debias_prog_dual}
		\min_{\ell \in \mathbb{R}^d} \left\{\frac{1}{4n} \sum_{i=1}^n \hat{\pi}_i \left[X_i^T \ell\right]^2 + x^T \ell +\frac{\gamma}{n}\norm{\ell}_1 \right\}.
	\end{equation}

\begin{assumption}[Sub-Gaussian covariate]
	\label{assump:sub_gaussian}
	The covariate vector $X\in \mathbb{R}^d$ is sub-Gaussian, \emph{i.e.}, $\norm{u^TX}_{\psi_2} \lesssim \left[\mathrm{E}\left(u^TXX^T u\right)\right]^{1/2}$ for any $u\in \mathbb{R}^d$.
\end{assumption}

\begin{assumption}[Lower eigenvalue bound on $\mathrm{E}\left(RXX^T\right)$]
	\label{assump:gram_lower_bnd}
	The minimum eigenvalue of the complete-case Gram matrix $\mathrm{E}\left(RXX^T\right) \in \mathbb{R}^{d\times d}$ is bounded away from zero, \emph{i.e.}, there exists a constant $\kappa_R>0$ such that 
	$$ \inf_{v\in \mathbb{S}^{d-1}} \mathrm{E}\left[R(X^Tv)^2\right] \geq \kappa_R^2.$$
\end{assumption}

\begin{assumption}[Stochastic control of the propensity score estimation]
	\label{assump:prop_consistent_rate}
	Given any $n\geq 1$ and $\delta \in (0,1)$, there exists $r_{\pi} \equiv r_{\pi}(n,\delta) >0$ such that $\mathrm{P}\left(\max_{1 \leq i \leq n}\left|\hat{\pi}_i -\pi_i \right| > r_{\pi}\right) <\delta$.
\end{assumption}

\begin{definition}[$r_{\ell}$-approximation to $\ell_0(x)$]
	\label{defn:popul_dual}
	For any $x\in \mathbb{R}^d$, we define an $r_{\ell}$-approximation $\tilde{\ell}(x)$ to the population dual solution $\ell_0(x)$ as:
	\begin{equation}
		\label{popul_dual_approx}
		\tilde{\ell}(x) \equiv \tilde{\ell}(x;r_{\ell}):= \argmin_{\ell(x)\in \mathbb{R}^d} \left\{\norm{\ell(x)}_0: \norm{\ell(x) - \ell_0(x)}_2 \leq r_{\ell}\norm{\ell_0(x)}_2 \right\},
	\end{equation}
	where $r_{\ell} \in \left[0,\frac{1}{2}\right]$ controls the accuracy of approximation.
\end{definition}

\begin{assumption}[Sparsity of $\tilde{\ell}(x)$]
	\label{assump:sparse_dual}
	The $r_{\ell}$-approximation $\tilde{\ell}(x) \in \mathbb{R}^d$ to the population dual solution $\ell_0(x)\in \mathbb{R}^d$ exists and satisfies
	$s_{\ell}(x) := \norm{\tilde{\ell}(x)}_0 \ll \min\{n,d\}$.
\end{assumption}

\begin{theorem}[Consistency of the dual solution $\hat{\ell}(x)$]
	\label{thm:dual_consist2}
	Let $\delta \in (0,1)$, $x\in \mathbb{R}^d$ be a fixed covariate vector, and $A_1,A_2,A_3>0$ be some large absolute constants. Suppose that Assumptions~\ref{assump:basic}, \ref{assump:sub_gaussian}, \ref{assump:gram_lower_bnd}, \ref{assump:prop_consistent_rate}, and \ref{assump:sparse_dual} hold
	as well as
	$$\sqrt{\frac{s_{\ell}(x)\log(ed/s_{\ell}(x))}{n}} + \sqrt{\frac{\log(1/\delta)}{n}} + r_{\pi} < \min\left\{\frac{\kappa_R^2}{A_1 \phi_{s_{\ell}(x)}}, 1 \right\} \quad \text{ and } \quad d\geq s_{\ell}(x) + 2.$$
	We denote $r_{\gamma}\equiv r_{\gamma}(n,\delta) = \frac{\sqrt{\phi_1}}{\kappa_R}\norm{x}_2\sqrt{\frac{\log(d/\delta)}{n}} + \frac{\sqrt{\phi_1 \phi_{s_{\ell}(x)}}\norm{x}_2}{\kappa_R^2} \cdot r_{\pi}$. If $\frac{\gamma}{n} = A_2\cdot r_{\gamma}$ and 
	$$r_{\ell} \leq \min\left\{\frac{1}{2},\, \left[\left(\frac{A_2^2 \,\kappa_R^6\, \phi_1}{A_3^2\norm{x}_2 \phi_{s_{\ell}(x)}}\right)^{\frac{1}{4}} \left(\frac{\log(d/\delta)}{n}\right)^{\frac{1}{4}} + \left(\frac{A_2^2 \phi_1 \kappa_R^4}{A_3^2 \phi_{s_{\ell}(x)} \norm{x}_2^2}\right)^{\frac{1}{4}} \sqrt{r_{\pi}}\right]\right\},$$
	then with probability at least $1-\delta$,
	\begin{align*}
		\norm{\hat{\ell}(x) - \tilde{\ell}(x)}_2
		&\lesssim \frac{\norm{x}_2}{\kappa_R^3}\left[ \sqrt{\frac{\phi_1 s_{\ell}(x)\log(d/\delta)}{n}} + \frac{\phi_{s_{\ell}(x)}}{\kappa_R} \cdot r_{\ell} + \frac{\sqrt{\phi_1 \phi_{s_{\ell}(x)} s_{\ell}(x)}}{\kappa_R} \cdot  r_{\pi} \right].
	\end{align*}
\end{theorem}

\begin{lemma}[Consistency of the dual solution $\hat{\ell}(x)$ when the propensity score is known]
	\label{lem:dual_consist}
	Let $\delta \in (0,1)$,  $x\in \mathbb{R}^d$ be a fixed covariate vector, and $A_1,A_2,A_3>0$ be some large absolute constants. Suppose that Assumptions~\ref{assump:basic}, \ref{assump:sub_gaussian}, \ref{assump:gram_lower_bnd}, and \ref{assump:sparse_dual} hold as well as
	$$\sqrt{\frac{s_{\ell}(x)\log(ed/s_{\ell}(x))}{n}} + \sqrt{\frac{\log(1/\delta)}{n}} < \min\left\{\frac{\kappa_R^2}{A_1\phi_{s_{\ell}(x)}}, 1 \right\} \quad \text{ and } \quad d\geq s_{\ell}(x)+2.$$
	If $\gamma = A_2 \frac{\sqrt{\phi_1}}{\kappa_R}\norm{x}_2 \sqrt{n\log(d/\delta)}$ and $r_{\ell} < \left(\frac{A_2^2 \,\kappa_R^6\, \phi_1}{A_3^2\norm{x}_2 \phi_{s_{\ell}(x)}}\right)^{\frac{1}{4}} \left(\frac{\log(d/\delta)}{n}\right)^{\frac{1}{4}}$, then with probability at least $1-\delta$,
	\begin{align*}
		\norm{\hat{\ell}(x)  - \tilde{\ell}(x)}_2 &\lesssim \frac{\norm{x}_2}{\kappa_R^2}\left[\sqrt{\frac{\phi_1}{\kappa_R^2}} \sqrt{\frac{s_{\ell}(x) \log(d/\delta)}{n}} + \frac{\phi_{s_{\ell}(x)}}{\kappa_R^2} \cdot r_{\ell} \right].
	\end{align*}
\end{lemma}

	\begin{align}
		\label{dual_taylor}
		\begin{split}
			&\left[\hat{\ell}(x) - \tilde{\ell}(x)\right]^T \left[\frac{1}{4n}\sum_{i=1}^n \pi(X_i) X_iX_i^T \right] \left[\hat{\ell}(x) - \tilde{\ell}(x)\right] \\
			&\quad + \left[\hat{\ell}(x) - \tilde{\ell}(x)\right]^T \left[\frac{1}{2n}\sum_{i=1}^n \pi(X_i) X_iX_i^T \right] \left[\tilde{\ell}(x) - \ell_0(x)\right]\\
			&\leq \left[\frac{1}{2n} \sum_{i=1}^n \pi(X_i) X_iX_i^T \ell_0(x) +x \right]^T \left[\tilde{\ell}(x) - \hat{\ell}(x) \right] + \frac{\gamma}{n}\left[\norm{\tilde{\ell}(x)}_1 - \norm{\hat{\ell}(x)}_1 \right].
		\end{split}
	\end{align}

\begin{proof}[Proof of \autoref{thm:dual_consist2}]
	Since $\hat{\ell}(x)$ minimizes the objective function of the dual problem \eqref{debias_prog_dual}, we know that
	$$\frac{1}{4n} \sum_{i=1}^n \hat{\pi}_i\left[X_i^T \hat{\ell}(x)\right]^2 + x^T\hat{\ell}(x) + \frac{\gamma}{n}\norm{\hat{\ell}(x)}_1 \leq \frac{1}{4n} \sum_{i=1}^n \hat{\pi}_i\left[X_i^T \tilde{\ell}(x)\right]^2 + x^T\tilde{\ell}(x) + \frac{\gamma}{n}\norm{\tilde{\ell}(x)}_1.$$
	Some rearrangements (or Taylor's expansion) give us that
	\begin{align}
		\label{dual_taylor2}
		\begin{split}
			&\left[\hat{\ell}(x) - \tilde{\ell}(x)\right]^T \left[\frac{1}{4n}\sum_{i=1}^n \hat{\pi}_i X_iX_i^T \right] \left[\hat{\ell}(x) - \tilde{\ell}(x)\right] + \left[\hat{\ell}(x) - \tilde{\ell}(x)\right]^T \left[\frac{1}{2n}\sum_{i=1}^n \hat{\pi}_i X_iX_i^T \right] \left[\tilde{\ell}(x) - \ell_0(x)\right]\\
			&\leq \left[\frac{1}{2n} \sum_{i=1}^n \hat{\pi}_i X_iX_i^T \ell_0(x) +x \right]^T \left[\tilde{\ell}(x) - \hat{\ell}(x) \right] + \frac{\gamma}{n}\left[\norm{\tilde{\ell}(x)}_1 - \norm{\hat{\ell}(x)}_1 \right].
		\end{split}
	\end{align}
	We now reduce the above inequality to \eqref{dual_taylor} in Lemma~\ref{lem:dual_consist} so that we can apply parts of the results in Lemma~\ref{lem:dual_consist}. Some algebra shows that the above inequality is equivalent to
	{\small\begin{align}
		\label{dual_taylor3}
		\begin{split}
			&\underbrace{\left[\hat{\ell}(x) - \tilde{\ell}(x)\right]^T \left[\frac{1}{4n}\sum_{i=1}^n \pi(X_i) X_iX_i^T \right] \left[\hat{\ell}(x) - \tilde{\ell}(x)\right] + \left[\hat{\ell}(x) - \tilde{\ell}(x)\right]^T \left[\frac{1}{2n}\sum_{i=1}^n \pi(X_i) X_iX_i^T \right] \left[\tilde{\ell}(x) - \ell_0(x)\right]}_{\textbf{Term I}}\\
			&\quad + \underbrace{\left[\hat{\ell}(x) - \tilde{\ell}(x)\right]^T \left[\frac{1}{4n}\sum_{i=1}^n \left(\hat{\pi}_i-\pi_i\right) X_iX_i^T \right] \left[\hat{\ell}(x) - \tilde{\ell}(x)\right]}_{\textbf{Term II}}\\
			&\quad + \underbrace{\left[\hat{\ell}(x) - \tilde{\ell}(x)\right]^T \left[\frac{1}{2n}\sum_{i=1}^n \left(\hat{\pi}_i - \pi_i\right) X_iX_i^T \right] \left[\tilde{\ell}(x) - \ell_0(x)\right]}_{\textbf{Term III}}\\
			&\leq \left[\frac{1}{2n} \sum_{i=1}^n \hat{\pi}_i X_iX_i^T \ell_0(x) +x \right]^T \left[\tilde{\ell}(x) - \hat{\ell}(x) \right] + \frac{\gamma}{n}\left[\norm{\tilde{\ell}(x)}_1 - \norm{\hat{\ell}(x)}_1 \right].
		\end{split}
	\end{align}}%
	We first obtain an upper bound on the right-hand side of the above inequality via the same arguments as in the proof of Lemma~\ref{lem:dual_consist}. We write $\hat{\Delta}(x) = \hat{\ell}(x) - \tilde{\ell}(x)$. By H{\"o}lder's inequality and the (reverse) triangle inequality, together with the sparsity of $\tilde{\ell}(x)$ under Assumption~\ref{assump:sparse_dual}, the right-hand side of \eqref{dual_taylor2} is upper bounded by
	\begin{align*}
		&\left[\frac{1}{2n} \sum_{i=1}^n \hat{\pi}_i X_iX_i^T \ell_0(x) +x \right]^T \left[\tilde{\ell}(x) - \hat{\ell}(x) \right] + \frac{\gamma}{n}\left[\norm{\tilde{\ell}(x)}_1 - \norm{\hat{\ell}(x)}_1 \right] \\
		&\leq \norm{\frac{1}{2n} \sum_{i=1}^n \hat{\pi}_i X_iX_i^T \ell_0(x) +x}_{\infty} \norm{\hat{\Delta}(x)}_1 \\
		&\quad + \frac{\gamma}{n} \left[\norm{\left[\tilde{\ell}(x)\right]_{S_{\ell}(x)}}_1 - \norm{\left[\hat{\Delta}(x)\right]_{S_{\ell}(x)} + \left[\tilde{\ell}(x)\right]_{S_{\ell}(x)}}_1 - \norm{\left[\hat{\ell}(x)\right]_{S_{\ell}(x)^c}}_1\right]\\
		&\leq \frac{\gamma}{5n} \left[\norm{\left[\hat{\Delta}(x)\right]_{S_{\ell}(x)}}_1 + \norm{\left[\hat{\Delta}(x)\right]_{S_{\ell}(x)^c}}_1 \right] + \frac{\gamma}{n} \left[\norm{\left[\hat{\Delta}(x) \right]_{S_{\ell}(x)}}_1 - \norm{\left[\hat{\Delta}(x) \right]_{S_{\ell}(x)^c}}_1 \right]\\
		&\leq \frac{6\gamma}{5n} \sqrt{s_{\ell}(x)} \norm{\hat{\Delta}(x)}_2
	\end{align*}
	where $\norm{\left[\hat{\Delta}(x) \right]_{S_{\ell}(x)}}_1 = \sum_{k\in S_{\ell}(x)} \left|\hat{\Delta}_k(x) \right|$.
	
	Next, we derive lower bounds for the terms on the left-hand side of \eqref{dual_taylor3}. We know from Lemma~\ref{lem:res_cross_poly} by $a_0=5$ that 
	$$\hat{\ell}(x)-\tilde{\ell}(x) \in \mathcal{C}_1\left(S_{\ell}(x), 3 \right) \bigcup B_d^{(1)}\left(0,\, \frac{13 r_{\ell}^2 \norm{x}_2^2}{2 \kappa_R^4} \cdot \frac{\nu}{\gamma}\right) \equiv \mathcal{C}_1\left(S_{\ell}(x), 3 \right) \bigcup B_d^{(1)}(0,r_1),$$ 
	where $B_d^{(1)}(0,r) = \left\{v\in \mathbb{R}^d: \norm{v}_1 \leq r\right\}$ and $r_1 \equiv \frac{13 r_{\ell}^2 \norm{x}_2^2}{2 \kappa_R^4} \cdot \frac{\nu}{\gamma}$.
	
	Now, by Lemma~\ref{lem:rel_poly_cost}, there exist absolute constants $A_2,A_2'>0$ such that 
	\begin{align*}
		&\norm{\frac{1}{n}\sum_{i=1}^n \hat{\pi}_i X_i X_i^T \ell_0(x) + 2x}_{\infty} \\
		&\leq A_2'\cdot \frac{\sqrt{\phi_1}}{\kappa_R}\norm{x}_2 \left[\sqrt{\frac{\log(d/\delta)}{n}} + \frac{\log(d/\delta)}{n} \right]\\
		&\quad + A_2'\cdot \frac{\sqrt{\phi_1 \phi_{s_{\ell}(x)}}\norm{x}_2}{\kappa_R^2} \left[1 + \sqrt{\frac{\log(d/\delta)}{n}} + \frac{\log(d/\delta)}{n}\right] r_{\pi}(n,\delta)\\
		&\lesssim A_2\cdot \frac{\sqrt{\phi_1}}{\kappa_R}\norm{x}_2\sqrt{\frac{\log(d/\delta)}{n}} + A_2\cdot \frac{\sqrt{\phi_1 \phi_{s_{\ell}(x)}}\norm{x}_2}{\kappa_R^2} \cdot r_{\pi}(n,\delta) \equiv \frac{\gamma}{n},
	\end{align*}
	where we apply our growth condition $\sqrt{\frac{s_{\ell}(x)\log(ed/s_{\ell}(x))}{n}} + \sqrt{\frac{\log(1/\delta)}{n}} < 1$. Similarly, we obtain from Lemma~\ref{lem:rel_poly_cost} that there exist absolute constants $A_3,A_3'>0$ such that with probability at least $1-\delta$,
	\begin{align*}
		&\sup_{v\in \mathcal{C}_2\left(S_{\ell}(x), \vartheta_{\ell}\right)\cap B_d(0,1)} \left|\frac{1}{n} \sum_{i=1}^n \hat{\pi}_i\left[X_i^Tv\right]^2 \right| \\
		&\leq \frac{2A_3'}{13}\cdot \phi_{s_{\ell}(x)}\left[1+ \sqrt{\frac{\log(1/\delta) + s_{\ell}(x) \log\left(d/s_{\ell}(x)\right)}{n}}\right] \left[1 + r_{\pi}(n,\delta)\right] \\
		&\leq \frac{2A_3}{13}\cdot \phi_{s_{\ell}(x)} \equiv \frac{\nu}{n},
	\end{align*}
	where we apply our growth condition $\sqrt{\frac{s_{\ell}(x)\log(ed/s_{\ell}(x))}{n}} + \sqrt{\frac{\log(1/\delta)}{n}} + r_{\pi} < 1$ again. Hence, with probability at least $1-\delta$,
	$$\frac{\kappa_R^2}{\norm{x}_2} \sqrt{\frac{2\gamma}{13\nu}} \geq \left(\frac{A_2^2 \,\kappa_R^6\, \phi_1}{A_3^2\norm{x}_2 \phi_{s_{\ell}(x)}}\right)^{\frac{1}{4}} \left(\frac{\log(d/\delta)}{n}\right)^{\frac{1}{4}} + \left(\frac{A_2^2 \phi_1 \kappa_R^4}{A_3^2 \phi_{s_{\ell}(x)} \norm{x}_2^2}\right)^{\frac{1}{4}} \sqrt{r_{\pi}} > r_{\ell},$$
	where the last inequality follows from our condition. In other words, with probability at least $1-\delta$, we have that $r_1= \frac{13 r_{\ell}^2 \norm{x}_2^2}{2 \kappa_R^4} \cdot \frac{\nu}{\gamma} < 1$, and 
	$$\left[\mathcal{C}_1\left(S_{\ell}(x), 3 \right) \bigcup B_d^{(1)}(0,r_1)\right] \bigcap \mathbb{S}^{d-1} = \mathcal{C}_1\left(S_{\ell}(x), 3 \right) \bigcap \mathbb{S}^{d-1}.$$
	
	$\bullet$ {\bf Term I:} This term coincides with the left-hand side of \eqref{dual_taylor} in the proof of Lemma~\ref{lem:dual_consist}. Thus, we conclude that there exist absolute constants $A,A_0>0$ such that with probability at least $1-\delta$, 
	\begin{align*}
		&\textbf{Term I} \\
		&\geq \left[\frac{\kappa_R^2}{4} - A\phi_{s_{\ell}(x)} \sqrt{\frac{\log(1/\delta)+s_{\ell}(x)\log(ed/s_{\ell}(x))}{n}} \right]\norm{\hat{\Delta}(x)}_2^2 \\
		&\quad - A\left[\sqrt{\frac{\log(1/\delta)+s_{\ell}(x)\log(ed/s_{\ell}(x))}{n}} + 1\right]\phi_{s_{\ell}(x)} \cdot r_{\ell} \norm{\ell_0(x)}_2\norm{\hat{\Delta}(x)}_2\\
		&\geq \left[\frac{\kappa_R^2}{4} - A_0\phi_{s_{\ell}(x)} \sqrt{\frac{\log(1/\delta)+s_{\ell}(x)\log(ed/s_{\ell}(x))}{n}} \right]\norm{\hat{\Delta}(x)}_2^2 - A_0 \phi_{s_{\ell}(x)} \cdot r_{\ell} \norm{\ell_0(x)}_2\norm{\hat{\Delta}(x)}_2,
	\end{align*}
	where we use our growth condition $\sqrt{\frac{s_{\ell}(x)\log(ed/s_{\ell}(x))}{n}} + \sqrt{\frac{\log(1/\delta)}{n}} < 1$ to obtain the second inequality.
	
	$\bullet$ {\bf Term II:} By Lemmas~\ref{lem:gram_mat_conc} and \ref{lem:cone_ball_bnd}, we have that with probability at least $1-\delta$,
	\begin{align*}
		\left|\textbf{Term II}\right| &\leq \norm{\hat{\Delta}(x)}_2^2 \sup_{u \in \mathcal{C}_1(S_{\ell}(x),3) \bigcap \mathbb{S}^{d-1}} \frac{1}{4n} \sum_{i=1}^n \left|\hat{\pi}_i - \pi_i\right| \left(X_i^T u\right)^2\\
		&\leq \norm{\hat{\Delta}(x)}_2^2 \sup_{u \in \mathcal{C}_1(S_{\ell}(x),3) \bigcap \mathbb{S}^{d-1}} \frac{r_{\pi}(n,\delta)}{4n} \sum_{i=1}^n  \left(X_i^T u\right)^2 \\
		&\lesssim \phi_{s_{\ell}(x)} \cdot r_{\pi}(n,\delta) \norm{\hat{\Delta}(x)}_2^2 \left[\sqrt{\frac{\log(1/\delta)+s_{\ell}(x)\log(ed/s_{\ell}(x))}{n}} + 1\right]\\
		&\lesssim \phi_{s_{\ell}(x)} \cdot r_{\pi}(n,\delta) \norm{\hat{\Delta}(x)}_2^2,
	\end{align*}
	where we again use our growth condition $\sqrt{\frac{s_{\ell}(x)\log(ed/s_{\ell}(x))}{n}} + \sqrt{\frac{\log(1/\delta)}{n}} < 1$.
	
	$\bullet$ {\bf Term III:} By Assumption~\ref{assump:sparse_dual}, Lemmas~\ref{lem:dual_approx_cone} and \ref{lem:cone_ball_bnd} with $\vartheta_{\ell}=\frac{r_{\ell}}{\sqrt{1-r_{\ell}^2}} < 1$ and $r_{\ell} \leq \frac{1}{2}$, 
	\begin{align*}
		&\left|\textbf{Term III}\right| \\
		&\leq \sup_{\substack{u \in \mathcal{C}_1(S_{\ell}(x),3) \bigcap \mathbb{S}^{d-1}\\v\in \mathcal{C}_2(S_{\ell}(x), 1) \cap \mathbb{S}^{d-1}}} \left[\frac{1}{2n} \sum_{i=1}^n \left|\hat{\pi}_i - \pi_i\right| \left|X_i^Tu \right| |X_i^Tv| \right] r_{\ell} \norm{\ell_0(x)}_2 \norm{\hat{\Delta}(x)}_2\\
		&\leq r_{\pi}(n,\delta) \sup_{\substack{u \in \mathcal{C}_1(S_{\ell}(x),3) \bigcap \mathbb{S}^{d-1}\\v\in \mathcal{C}_2(S_{\ell}(x), 1) \cap \mathbb{S}^{d-1}}} \left[\frac{1}{2n} \sum_{i=1}^n \left|X_i^Tu \right| |X_i^Tv| \right] r_{\ell} \norm{\ell_0(x)}_2 \norm{\hat{\Delta}(x)}_2\\
		&\leq r_{\pi}(n,\delta) \cdot r_{\ell} \norm{\ell_0(x)}_2 \norm{\hat{\Delta}(x)}_2 \cdot \phi_{s_{\ell}(x)} \left[\sqrt{\frac{\log(1/\delta) + s_{\ell}(x)\log(ed/s_{\ell}(x))}{n}} + 1 \right]
	\end{align*}
	Under the growth condition $\sqrt{\frac{s_{\ell}(x)\log(ed/s_{\ell}(x))}{n}} + \sqrt{\frac{\log(1/\delta)}{n}} < 1$, we have that with probability at least $1-\delta$, 
	$$\left|\textbf{Term III}\right| \lesssim \phi_{s_{\ell}(x)} \cdot r_{\pi}(n,\delta) \cdot r_{\ell} \norm{\ell_0(x)}_2 \norm{\hat{\Delta}(x)}_2.$$
	
	Combining \eqref{dual_taylor3} with the above three bounds yields that there exists an absolute constant $A_1>0$ such that with probability at least $1-\delta$,
	\begin{align*}
		&\left[\frac{\kappa_R^2}{4} - \frac{A_1}{16}\phi_{s_{\ell}(x)} \sqrt{\frac{\log(1/\delta)+s_{\ell}(x)\log(ed/s_{\ell}(x))}{n}} -\frac{A_1}{16} \cdot \phi_{s_{\ell}(x)} \cdot r_{\pi}(n,\delta) \right]\norm{\hat{\Delta}(x)}_2^2 \\
		&\quad - A_1 \left[1+r_{\pi}(n,\delta)\right] \phi_{s_{\ell}(x)} r_{\ell} \norm{\ell_0(x)}_2\norm{\hat{\Delta}(x)}_2\\
		&\leq \frac{6\gamma}{5n} \sqrt{s_{\ell}(x)} \norm{\hat{\Delta}(x)}_2.
	\end{align*}
	By our growth condition $\sqrt{\frac{s_{\ell}(x)\log(ed/s_{\ell}(x))}{n}} + \sqrt{\frac{\log(1/\delta)}{n}} + r_{\pi} < \min\left\{\frac{\kappa_R^2}{A_1 \phi_{s_{\ell}(x)}}, 1 \right\}$, we conclude by dividing both sides of the inequality by $\norm{\hat{\Delta}(x)}_2$ and $\frac{\kappa_R^2}{8}$ that
	\begin{align*}
		\norm{\hat{\ell}(x) - \tilde{\ell}(x)}_2 &\leq \frac{48\gamma}{5\kappa_R^2 n} \sqrt{s_{\ell}(x)} + \frac{8A_1}{\kappa_R^2} \left[1+r_{\pi}(n,\delta)\right] \phi_{s_{\ell}(x)} r_{\ell} \norm{\ell_0(x)}_2\\
		&\lesssim \frac{\gamma \sqrt{s_{\ell}(x)}}{n \cdot \kappa_R^2} + \frac{\phi_{s_{\ell}(x)} \norm{x}_2}{\kappa_R^4} \cdot r_{\ell}.
	\end{align*}
	Finally, the bound for $\norm{\hat{\ell}(x)-\tilde{\ell}(x)}_2$ follows by substituting the asymptotic rate of $\frac{\gamma}{n}$ into the upper bound of $\norm{\hat{\ell}(x)-\tilde{\ell}(x)}_2$.
\end{proof}

\begin{lemma}
	\label{lem:gram_mat_conc}
	Let $\vartheta \in [0,\infty]$ and $\delta \in (0,1)$ be arbitrary. Suppose that Assumption~\ref{assump:sub_gaussian} holds. Let $S\subset \{1,...,d\}$ with $|S|=s$. If $d\geq s+2$, then 
	$$\mathrm{E}\left[\sup_{v\in \mathcal{C}_1(S,\vartheta) \cap B_d(0,1)} \left|\frac{1}{n}\sum_{i=1}^n \Big\{(X_i^T v)^2 - \mathrm{E}\left[(X_i^T v)^2 \right]\Big\}\right|\right] \lesssim (2+\vartheta)^2 \phi_s \left[\sqrt{\frac{s\log(ed/s)}{n}} + \frac{s\log(ed/s)}{n} \right]$$
	and with probability at least $1-\delta$,
	\begin{align*}
		&\sup_{v\in \mathcal{C}_1(S,\vartheta) \cap B_d(0,1)} \left|\frac{1}{n}\sum_{i=1}^n \Big\{(X_i^T v)^2 - \mathrm{E}\left[(X_i^T v)^2 \right]\Big\}\right| \\
		&\lesssim (2+\vartheta)^2 \phi_s \left[\sqrt{\frac{\log(1/\delta) + s\log(ed/s)}{n}} + \frac{\log(1/\delta)  + s\log(ed/s)}{n} \right].
	\end{align*}
	In addition, these bounds also apply to
	$$\sup_{v\in \mathcal{C}_1(S,\vartheta) \cap B_d(0,1)} \left|\frac{1}{n}\sum_{i=1}^n \Big\{R_i(X_i^T v)^2 - \mathrm{E}\left[R_i(X_i^T v)^2 \right]\Big\}\right|$$
	for the complete-case analysis and 
	$$\sup_{v\in \mathcal{C}_1(S,\vartheta) \cap B_d(0,1)} \left|\frac{1}{n}\sum_{i=1}^n \left\{\frac{R_i}{\pi_i}(X_i^T v)^2 - \mathrm{E}\left[\frac{R_i}{\pi_i}(X_i^T v)^2 \right]\right\}\right|$$
	for the IPW analysis with the true propensity scores $\pi_i = \mathrm{P}(R_i=1|X_i), i=1,...,n$ under Assumption~\ref{assump:basic}.
	Finally, the above asymptotic upper bound remains almost identical for 
	\begin{align*}
		&\sup_{\substack{u \in \mathcal{C}_2(S_1,\vartheta_2) \cap \mathbb{S}^{d-1}\\v\in \mathcal{C}_1(S_2,\vartheta_1) \cap \mathbb{S}^{d-1}}} \left|\frac{1}{n}\sum_{i=1}^n \Big\{(X_i^T u)(X_i^T v) - \mathrm{E}\left[(X_i^Tu)(X_i^T v) \right]\Big\}\right|\\
		&\lesssim (2+\vartheta_1) (2+\vartheta_2) \sqrt{\phi_{s_1}\phi_{s_2}}\\
		&\quad \times  \left[\sqrt{\frac{\log(1/\delta) + s_1\log(ed/s_1) + s_2\log(ed/s_2)}{n}} + \frac{\log(1/\delta) + s_1\log(ed/s_1) + s_2\log(ed/s_2)}{n} \right]
	\end{align*}
	with $|S_1|=s_1$ and $|S_2|=s_2$.
\end{lemma}

\begin{lemma}
	\label{lem:dual_approx_cone}
	Under Definition~\ref{defn:popul_dual} and Assumption~\ref{assump:sparse_dual}, we have that
	$$\ell_0(x) \in \mathcal{C}_2\left(S_{\ell}(x),\vartheta_{\ell}\right)=\left\{v\in \mathbb{R}^d: \norm{v_{S_{\ell}(x)^c}}_2 \leq \vartheta_{\ell} \norm{v_{S_{\ell}(x)}}_2 \right\}\quad \text{ with } \quad \vartheta_{\ell}=\frac{r_{\ell}}{\sqrt{1-r_{\ell}^2}}$$
	for any $r_{\ell}\in [0,1)$. Moreover, if $0\leq r_{\ell} \leq \frac{1}{2}$, then $$\norm{M^T\left(\tilde{\ell}(x)-\ell_0(x)\right)}_2 \leq \sup_{v\in \mathcal{C}_2(S_{\ell}(x),1) \cap B_d\left(0,10r_{\ell}\norm{\ell_0(x)}_2\right)} \norm{M^T v}_2$$ 
	for any matrix $M\in \mathbb{R}^{d\times k}$.
\end{lemma}

\begin{lemma}[Restricted cone and cross polytope property]
	\label{lem:res_cross_poly}
	Let $a_0 >1$ be some absolute constant. Under Assumptions~\ref{assump:gram_lower_bnd} and \ref{assump:sparse_dual}, if 
	\begin{equation}
		\label{res_cross_poly_cond}
		\norm{\frac{1}{2n}\sum_{i=1}^n \hat{\pi}_i\cdot \ell_0(x)^T X_i X_i +x}_{\infty} \leq \frac{\gamma}{a_0 n} \quad \text{ and } \quad \sup_{v\in \mathcal{C}_2\left(S_{\ell}(x), \vartheta_{\ell}\right)\cap B_d(0,1)} \left|\frac{1}{n} \sum_{i=1}^n \hat{\pi}_i\left[X_i^Tv\right]^2 \right| \leq \frac{\nu}{a_0 n}
	\end{equation}
	with $\vartheta_{\ell} = \frac{r_{\ell}}{\sqrt{1-r_{\ell}^2}}$, then
	$$\sum_{k\in S_{\ell}(x)^c} \left|\hat{\ell}_k(x) \right| \leq \frac{4r_{\ell}^2 \norm{x}_2^2}{(a_0-1) \kappa_R^4}\cdot \frac{\nu}{\gamma} + \left(\frac{a_0+1}{a_0-1}\right) \sum_{k\in S_{\ell}(x)^c} \left|\hat{\ell}_k(x) -\tilde{\ell}_k(x) \right|$$
	so that 
	$$\hat{\ell}(x)-\tilde{\ell}(x) \in \mathcal{C}_1\left(S_{\ell}(x), 2\left(\frac{a_0+1}{a_0-1}\right) \right) \bigcup B_d^{(1)}\left(0,\, \frac{(a_0^2+1)r_{\ell}^2 \norm{x}_2^2}{(a_0-1) \kappa_R^4} \cdot \frac{\nu}{\gamma}\right),$$ where $B_d^{(1)}(0,r) = \left\{v\in \mathbb{R}^d: \norm{v}_1 \leq r\right\}$.
\end{lemma}

\begin{lemma}
	\label{lem:cone_ball_bnd}
	Let $\vartheta,r_0,r_1 \in (0,\infty)$ and $\delta \in (0,1)$ be arbitrary. Suppose that Assumption~\ref{assump:sub_gaussian} holds. Let $S,\tilde{S}\subset \{1,...,d\}$ with $|S|=s$ and $|\tilde{S}|=\tilde{s}$. If $d\geq s+2$, then with probability at least $1-\delta$,
	\begin{align*}
		&\sup_{v\in \left[\mathcal{C}_1(S,\vartheta) \bigcup B_d^{(1)}(0,r_1) \right] \bigcap B_d(0,r_0)} \left|\frac{1}{n} \sum_{i=1}^n \left\{(X_i^Tv)^2 - \mathrm{E}\left[(X_i^Tv)^2\right] \right\}\right| \\
		&\lesssim r_0^2 (2+\vartheta)^2 \phi_s\left[\sqrt{\frac{\log(1/\delta) + s\log(ed/s)}{n}} + \frac{\log(1/\delta)  + s\log(ed/s)}{n} \right] \\
		&\quad + r_1^2\phi_1 \left[\sqrt{\frac{\log(d/\delta)}{n}} + \frac{\log(d/\delta)}{n}\right].
	\end{align*}
	In particular, when $r_0\geq r_1$, we further have that with probability at least $1-\delta$,
	\begin{align*}
		&\sup_{v\in \left[\mathcal{C}_1(S,\vartheta) \bigcup B_d^{(1)}(0,r_1) \right] \bigcap B_d(0,r_0)} \left|\frac{1}{n} \sum_{i=1}^n \left\{(X_i^Tv)^2 - \mathrm{E}\left[(X_i^Tv)^2\right] \right\}\right| \\
		&\lesssim r_0^2 (2+\vartheta)^2 \phi_s\left[\sqrt{\frac{\log(1/\delta) + s\log(ed/s)}{n}} + \frac{\log(1/\delta)  + s\log(ed/s)}{n} \right].
	\end{align*}
	Similarly, if $d\geq \max\{s,\tilde{s}\} +2$, then with probability at least $1-\delta$,
	\begin{align*}
		&\sup_{\substack{u \in \mathcal{C}_2(\tilde{S},\vartheta) \cap B_d(0,1)\\v\in \left[\mathcal{C}_1(S,\vartheta) \bigcup B_d^{(1)}(0,r_1) \right] \bigcap B_d(0,r_0)}} \frac{1}{n} \sum_{i=1}^n \left|X_i^Tu \right| \left|X_i^Tv \right| \\
		&\lesssim \left[r_0(2+\vartheta)^2+r_1\right] \sqrt{\phi_{\tilde{s}} \cdot \phi_1} \left[\sqrt{\frac{\log(1/\delta) + s\log(ed/s)}{n}} + \frac{\log(1/\delta)  + s\log(ed/s)}{n} + 1\right].
	\end{align*}
\end{lemma}

\begin{lemma}
	\label{lem:rel_poly_cost}
	Let $\delta\in (0,1)$ and $x\in \mathbb{R}^d$ be a fixed covariate vector. Suppose that Assumption~\ref{assump:sub_gaussian} holds. Then, with probability at least $1-\delta$,
	\begin{align*}
		&\sup_{v\in \mathcal{C}_2\left(S_{\ell}(x), \vartheta_{\ell}\right)\cap B_d(0,1)} \left|\frac{1}{n} \sum_{i=1}^n \pi_i\left[X_i^Tv\right]^2 \right| \\
		&\lesssim (2+\vartheta_{\ell})^2 \phi_{s_{\ell}(x)}\left[1+\sqrt{\frac{\log(\frac{1}{\delta})+s_{\ell}(x)\log\left(\frac{ed}{s_{\ell}(x)}\right)}{n}} + \frac{\log(\frac{1}{\delta})+s_{\ell}(x)\log\left(\frac{ed}{s_{\ell}(x)}\right)}{n}\right].
	\end{align*}
	If, in addition, Assumption~\ref{assump:prop_consistent_rate} holds, then with probability at least $1-\delta$,
	\begin{align*}
		&\sup_{v\in \mathcal{C}_2\left(S_{\ell}(x), \vartheta_{\ell}\right)\cap B_d(0,1)} \left|\frac{1}{n} \sum_{i=1}^n \hat{\pi}_i\left[X_i^Tv\right]^2 \right| \\
		&\lesssim \left[1+r_{\pi}(n,\delta)\right] (2+\vartheta_{\ell})^2 \phi_{s_{\ell}(x)}\left[1+\sqrt{\frac{\log(\frac{1}{\delta})+s_{\ell}(x)\log\left(\frac{ed}{s_{\ell}(x)}\right)}{n}} + \frac{\log(\frac{1}{\delta})+s_{\ell}(x)\log\left(\frac{ed}{s_{\ell}(x)}\right)}{n}\right].
	\end{align*}
\end{lemma}

\end{document}
